\documentclass[11pt]{article}
\usepackage[margin=1in]{geometry}
\usepackage{booktabs}
\usepackage{hyperref}
\usepackage{xcolor}
\usepackage{enumitem}
\usepackage{fancyvrb}
\hypersetup{colorlinks=true, linkcolor=blue, urlcolor=blue}

\DefineVerbatimEnvironment{code}{Verbatim}{fontsize=\small,frame=single,framesep=4pt,rulecolor=\color{black!20}}

\title{Reproducibility Report \\ \large Project~2: Longitudinal CEA Trajectories and Disease Progression Prediction in Colorectal Cancer}
\author{Yanan Fang}
\date{April 2026}

\begin{document}
\maketitle

This report documents the reproducibility practices used for Project~2. It is organized around the five practice areas covered in the course's reproducibility lecture: file organization and naming, version control, safe programming, dynamic reports, and pipelines.

\paragraph{Project links.}
\begin{itemize}[leftmargin=*]
  \item Code repository (GitHub): \url{https://github.com/yananf8899/Biostat629}
  \item Project website (GitHub Pages): \url{https://yananf8899.github.io/Biostat629/}
  \item Manuscript PDF (in repository): \texttt{project2\_report.pdf}
\end{itemize}

\section{File Organization and Naming}

The project follows a flat layout with descriptive, lowercase, underscore-separated file names. The lecture's recommendation to date-stamp working files is applied to interim notes; the final deliverables use stable names so that external links (e.g., the website's \texttt{report.pdf}) do not break across edits.

\begin{code}
.
|-- README.md                   project-level reproducibility notes
|-- Project2_Proposal.md        original proposal
|-- project2_analysis.R         single-script analysis pipeline
|-- project2_report.tex/.pdf    final manuscript
|-- project2_presentation.pptx  in-class slides
|-- project2_poster.pptx        poster
|-- reproducibility_report.tex/.pdf  this document
|-- sessionInfo.txt             package versions captured at run time
|-- fig{1..5}_*.pdf / .png      figures generated by the R script
|-- docs/                       GitHub Pages website
|   |-- index.html
|   |-- style.css
|   |-- figures/
|   `-- report.pdf
`-- msk_chord_2024/             controlled-access data (NOT in git)
\end{code}

The raw input directory \texttt{msk\_chord\_2024/} is referenced by a single relative path variable inside the R script (\texttt{data\_dir <- "msk\_chord\_2024"}), and the directory itself is excluded from version control via \texttt{.gitignore} because MSK-CHORD is controlled-access data that may not be redistributed.

\section{Version Control}

The project is tracked with \texttt{git} and hosted on GitHub at \url{https://github.com/yananf8899/Biostat629}. The repository contains every file needed to reproduce the manuscript, the website, and the figures, except the controlled-access input data. The git history records the development state of the analysis and manuscript and allows the project to be versioned and shared transparently.

A minimal \texttt{.gitignore} excludes the raw data folder and intermediate LaTeX build artifacts:

\begin{code}
msk_chord_2024/
*.aux
*.log
*.out
.DS_Store
\end{code}

\section{Safe Programming Practices}

The analysis script \texttt{project2\_analysis.R} follows the practices recommended in lecture:

\begin{itemize}[leftmargin=*]
  \item \textbf{Raw data is never modified in place.} All filtering, reshaping, and feature construction happens in R memory using \texttt{data.table}; the files in \texttt{msk\_chord\_2024/} are read but never written.
  \item \textbf{No use of \texttt{setwd()}, \texttt{rm()}, \texttt{source()}, or \texttt{system()}.} The data path is a single top-of-file variable (\texttt{data\_dir}); the script does not change the working directory or call out to the shell.
  \item \textbf{No interactive \texttt{save.image()} / \texttt{load(".RData")}.} Each run starts from a clean session.
  \item \textbf{Random seeds are fixed.} \texttt{set.seed(629)} is called immediately before each random step: the patient-level 70/30 train/test split, the LASSO 5-fold cross-validation foldid, the conformal calibration sub-split, and the figure-1 patient sample. Every reported number reproduces exactly.
  \item \textbf{No hard-coded results.} All numbers cited in the manuscript and on the website (AUCs, $p$-values, calibration slopes, conformal coverage, death-in-window rates) are printed by the script; no value is typed by hand.
  \item \textbf{Numerical guards.} CEA values are clamped to non-negative before \texttt{log()} via the offset \texttt{log(RESULT + 0.1)}; predicted probabilities are clamped to $[10^{-6}, 1-10^{-6}]$ before being passed to \texttt{qlogis()} for the calibration fit.
  \item \textbf{Patient-level integrity.} Train/test and CV folds are formed at the \emph{patient} level (not the landmark level), so no patient appears in both train and test, and the LASSO CV foldid uses a per-patient lookup so all of a patient's landmarks share a fold.
  \item \textbf{Diagnostic and audit blocks.} The script prints (i) cohort attrition counts, (ii) the death-in-window competing-risk rate, and (iii) a high-precision Model~A vs.\ Model~B AUC audit (six-decimal AUCs and rank correlation), so reviewers can verify the headline numbers without re-running ad hoc analyses.
\end{itemize}

\section{Dynamic Reports / Environment Capture}

The manuscript is written in plain \LaTeX{} (\texttt{project2\_report.tex}) rather than R~Markdown. The lecture recommends dynamic reports because they tie text and code together, but the project uses a workable equivalent: a single deterministic R script writes every figure to disk under a fixed filename, and the \LaTeX{} document includes those filenames. Re-running the R script and recompiling \texttt{pdflatex} regenerates the manuscript end-to-end.

To capture the software environment, the script ends with:

\begin{code}
capture.output(sessionInfo(), file = "sessionInfo.txt")
\end{code}

\noindent so that R's version, the loaded packages, and their versions for the most recent run are committed to the repository as \texttt{sessionInfo.txt}. Anyone re-running the analysis can compare against this file to detect environment drift.

\section{Reproducible Pipeline}

The project does not require a Snakemake or Docker pipeline because the analysis is small enough to fit in a single R script. The pipeline is deliberately one command:

\begin{code}
Rscript project2_analysis.R          # generates fig1..fig5 and sessionInfo.txt
pdflatex project2_report.tex         # compiles the manuscript (twice for refs)
pdflatex project2_report.tex
pdflatex reproducibility_report.tex  # compiles this document
\end{code}

Running these four commands on a fresh checkout (with \texttt{msk\_chord\_2024/} placed at the project root) regenerates every figure, every table, every number cited in the manuscript and on the website, and the manuscript and reproducibility-report PDFs. Run time is approximately three minutes on a 2023 MacBook Pro, dominated by the \texttt{glmer()} fit for Model~A.

\paragraph{Software requirements.}
\begin{itemize}[leftmargin=*]
  \item R~$\geq 4.3.0$
  \item R packages: \texttt{data.table}, \texttt{lme4}, \texttt{pROC}, \texttt{ggplot2}, \texttt{glmnet}
  \item TeX~Live with \texttt{pdflatex}
\end{itemize}

\noindent Exact versions are captured in \texttt{sessionInfo.txt}.

\section{Data Access}

MSK-CHORD~2024 is controlled-access. Reviewers must obtain the release through cBioPortal (\url{https://cbioportal.org}) and unpack it into a folder named \texttt{msk\_chord\_2024/} at the project root. The script expects the following seven files inside that folder:

\begin{code}
data_clinical_patient.txt
data_clinical_sample.txt
data_timeline_cea_labs.txt
data_timeline_progression.txt
data_timeline_diagnosis.txt
data_timeline_treatment.txt
data_mutations.txt
\end{code}

\noindent No data are committed to the repository, and the script does not modify any of these files.

\section{Summary Checklist}

\begin{tabular}{ll}
\toprule
Practice & Status \\
\midrule
Descriptive file names, flat directory layout & yes \\
Project tracked in \texttt{git}, hosted on GitHub & yes (link above) \\
\texttt{.gitignore} excludes controlled-access data and build artifacts & yes \\
Raw data never modified in place & yes \\
No \texttt{setwd}, \texttt{rm}, \texttt{source}, \texttt{system} in scripts & yes \\
Random seeds fixed for every stochastic step & yes (\texttt{set.seed(629)}) \\
Patient-level train/test and CV folds (no leakage) & yes \\
All cited numbers printed by the script (no hand-typed values) & yes \\
Software environment captured in \texttt{sessionInfo.txt} & yes \\
Single-command reproduction & yes (4-line shell pipeline) \\
Project website with link to repo and report & yes (link above) \\
\bottomrule
\end{tabular}

\end{document}
